\section{Experimental Setup}
\label{sec:setup}

\paragraph{Tasks and models.}
We use three configurations built on BizBench~\cite{bizbench2024}, chosen so
that the dataset family, the model family, the backbone scale, and the strength
of the deterministic baseline all vary:

\begin{itemize}\itemsep2pt
\item \textbf{Config~1}: CodeTAT-QA $\times$ Qwen2.5-0.5B-Instruct. Context is a
JSON table; gold programs index cells as \texttt{df["col"]["row"]}. 2846 clean
training items, 400 held out as dev (all selection uses a fixed 200-item subset
of it), frozen test $n{=}282$.
\item \textbf{Config~2}: CodeFinQA $\times$ Llama-3.2-1B-Instruct. Context is
narrative financial-report text with embedded tables; gold programs inline the
numbers. 4035 clean training items, 400 dev (same 200-item convention), frozen
test $n{=}664$.
\item \textbf{Config~3}: CodeTAT-QA $\times$ Qwen2.5-7B-Instruct. Identical
task, data splits, budgets, and protocol to Config~1 with a $14\times$ larger
backbone; this is the scale probe of Section~\ref{sec:scale}.
\end{itemize}

An item is \emph{clean} if its gold program executes to its gold answer under
our verifier: $2846/2856 = 99.65\%$ for CodeTAT-QA and $4035/4668 = 86.4\%$ for
CodeFinQA. The ten CodeTAT-QA failures were inspected individually and are
dataset defects (a final line not assigned to \texttt{answer}, or a
comprehension with no bare binary operation), not adaptation artifacts.

\paragraph{Demonstration budgets.}
Nested subsets are drawn from the clean pool, so a smaller budget is always a
prefix of a larger one. Config~1 uses
$m \in \{0,1,2,4,8,16,64,256,1024,2446\}$ and Config~2 uses
$m \in \{0,16,64,128,256,1024,3635\}$, where the largest value is the full
pool. Three independent subset seeds are drawn per configuration; every arm at
a given $(m, \text{seed})$ sees exactly the same $m$ demonstrations.

\paragraph{Pre-registration.}
The decision quantity in Eq.~\eqref{eq:delta}, the gate criteria, the budget
selection rule, and the verdict vocabulary were frozen before any training or
evaluation run on the second task. The protocol also fixed, in advance, how a
result contrary to our expectation would be handled: \emph{if the solver is
strong and GRPO nevertheless wins, report it as such, and do not weaken the
solver to recover the original narrative}. Both outcomes were accepted at
freeze time. The solver's construction was likewise constrained ex ante to a
``good-faith maximum effort'' clause forbidding template pruning or threshold
tuning in response to intermediate results.

\paragraph{Budget selection.}
Two operating points are chosen per configuration from the warm-start dose
curve using the pre-registered bands: $m_L$ is the smallest budget landing in
$[5\%, 20\%]$ dev accuracy and $m_H$ the smallest in $[25\%, 45\%]$. This
yields $m_L{=}2, m_H{=}64$ for Config~1 and $m_L{=}128, m_H{=}256$ for
Config~2. When the initial grid left the low band empty, additional budgets
were run rather than reinterpreting the band---this is one of the logged
deviations.

\paragraph{Leakage control.}
We assert disjointness of dev, demonstration pool, and test, and verify before
each solver run that its retrieval library does not intersect dev. The frozen
test set is untouched during training, checkpoint selection, and budget
selection; every arm is evaluated on it exactly once, at the checkpoint already
selected on dev.

\paragraph{Deviations.}
Minor protocol adjustments (e.g., adding intermediate budget points when initial grids left required bands empty) were logged alongside their influence direction.
